\subsection{Zeilenstufenform, Pivotpositionen und Rang}
\label{subsec:linear_systems_row_echelon_rank}

Ziel der Vorwärtselimination ist nicht nur eine Dreiecksform für quadratische
Matrizen, sondern die allgemeinere \emph{Zeilenstufenform}.

\begin{DefinitionBox}{Zeilenstufenform}{linear_systems_row_echelon_form}
  Eine Matrix besitzt Zeilenstufenform, wenn
  \begin{enumerate}[label=\arabic*.,itemsep=0.25em]
    \item alle Nullzeilen unterhalb der Nichtnullzeilen stehen,
    \item der erste von null verschiedene Eintrag jeder Nichtnullzeile weiter
      rechts steht als derjenige der vorherigen Zeile und
    \item unter jedem solchen führenden Eintrag nur Nullen stehen.
  \end{enumerate}
\end{DefinitionBox}

\begin{DefinitionBox}{Pivotelement und Pivotvariable}{linear_systems_pivot_terms}
  In einer Zeilenstufenform heißt der erste von null verschiedene Eintrag
  jeder Nichtnullzeile \emph{Pivotelement}. Seine Stelle in der Matrix heißt
  \emph{Pivotposition}; die zugehörige Spalte der Koeffizientenmatrix heißt
  \emph{Pivotspalte}. Die Unbekannte einer Pivotspalte ist eine
  \emph{Pivotvariable}.

  Eine Variablenspalte ohne Pivotposition gehört zu einer \emph{freien
  Variablen}. Ein Pivotelement muss nicht den Wert $1$ besitzen.
\end{DefinitionBox}

\begin{center}
  \begin{tikzpicture}[
  >=Latex,
  every node/.style={font=\small},
  pivot/.style={draw=THMBlue,line width=0.9pt,rounded corners=2pt,inner sep=2pt},
  guide/.style={THMBlue,line width=1pt},
  brace/.style={decorate,decoration={brace,amplitude=4pt}}
]
  \matrix (M) [
    matrix of math nodes,
    left delimiter={[},right delimiter={]},
    nodes={minimum width=0.95cm,minimum height=0.68cm,anchor=center},
    column sep=1pt,row sep=1pt
  ] {
    \star&\ast&\ast&\ast&\ast&\ast\\
    0&\star&\ast&\ast&\ast&\ast\\
    0&0&0&\star&\ast&\ast\\
    0&0&0&0&0&0\\
  };

  \coordinate (bar-top) at ($(M-1-5.north east)!0.5!(M-1-6.north west)+(0,0.18)$);
  \coordinate (bar-bottom) at ($(M-4-5.south east)!0.5!(M-4-6.south west)-(0,0.18)$);
  \draw[line width=0.65pt] (bar-top)--(bar-bottom);

  \node[pivot,fit=(M-1-1)] (p1) {};
  \node[pivot,fit=(M-2-2)] (p2) {};
  \node[pivot,fit=(M-3-4)] (p3) {};
  \draw[guide,densely dotted] (p1.south east)--(p2.north west);
  \draw[guide,densely dotted] (p2.south east)--(p3.north west);

  \node[text=THMBlue,align=left,anchor=west] (pivots)
    at ($(M.east)+(1.25,0.9)$)
    {Pivotelemente\\bilden eine Treppe};
  \draw[->,THMBlue] (pivots.west)--(p3.east);

  \draw[brace,black!55]
    ($(M-4-1.south west)+(0,-0.22)$)--($(M-4-5.south east)+(0,-0.22)$)
    node[midway,below=6pt,text=black!65]
    {$n$ Variablenspalten};
  \node[anchor=west,text=black!65,align=left]
    at ($(M.east)+(1.25,-0.15)$)
    {hier: $3$ Pivotpositionen\\$\Rightarrow\operatorname{rang}(A)=3$};
  \node[text=black!55,anchor=west] (zero)
    at ($(M.east)+(1.25,-1.05)$) {Nullzeile};
  \draw[->,black!50] (zero.west)--(M-4-6.east);
\end{tikzpicture}

\end{center}

Freie Variablen können bei einem lösbaren LGS als Parameter gewählt werden.

\clearpage
\begin{samepage}
  \begin{DefinitionBox}{Rang einer Matrix}{linear_systems_matrix_rank}
    Der \emph{Rang} einer Matrix ist die Anzahl ihrer Pivotpositionen in einer
    Zeilenstufenform. Gleichbedeutend ist er die Anzahl der Nichtnullzeilen
    einer Zeilenstufenform. Man schreibt $\operatorname{rang}(A)$.
  \end{DefinitionBox}

  Elementare Zeilenumformungen verändern den Rang nicht. Obwohl eine Matrix
  auf verschiedenen Wegen in Zeilenstufenform gebracht werden kann, ist die
  Anzahl ihrer Pivotpositionen daher eindeutig.

  \begin{ExampleBox}{Rang und freie Variablen ablesen}{linear_systems_rank_example}
    \[
      A\sim
      \left[
        \begin{array}{rrrr}
          1&2&0&-1\\
          0&0&1& 4\\
          0&0&0& 0
        \end{array}
      \right].
    \]
    Die Pivotpositionen liegen in Spalte $1$ und $3$. Somit gilt
    \[
      \operatorname{rang}(A)=2,
      \qquad
      x_1,x_3\text{ sind Pivotvariablen},
      \qquad
      x_2,x_4\text{ sind frei}.
    \]
    Die Rangzahl ist nicht die Zahl der Spalten, sondern die Zahl der
    voneinander unabhängigen Stufen.
  \end{ExampleBox}
\end{samepage}

Das Gauß-Verfahren verbindet diese Begriffe zu einem festen Ablauf:

\begin{center}
  \begin{tikzpicture}[
  >=Latex,
  processbox/.style={
    draw=black!25,fill=black!2,rounded corners=2pt,
    minimum height=1.35cm,text width=3cm,align=center,font=\footnotesize
  },
  arrow/.style={->,THMBlue,line width=0.9pt}
]
  \node[processbox] (system) {LGS\\$Ax=b$};
  \node[processbox,right=0.3cm of system] (matrix)
    {erweiterte Matrix\\$[A\mid b]$};
  \node[processbox,right=0.3cm of matrix] (forward)
    {Vorwärts-\\elimination\\Zeilenstufenform};
  \node[processbox,right=0.3cm of forward] (decision)
    {Rang und\\Lösbarkeit prüfen};
  \node[processbox,right=0.3cm of decision] (back)
    {Rückwärts-\\einsetzen oder\\parametrisieren};

  \draw[arrow] (system)--(matrix);
  \draw[arrow] (matrix)--(forward);
  \draw[arrow] (forward)--(decision);
  \draw[arrow] (decision)--(back);
\end{tikzpicture}

\end{center}

\begin{KeyPointBox}{Gauß-Verfahren}{linear_systems_gaussian_algorithm}
  \begin{enumerate}[label=\arabic*.,itemsep=0.3em]
    \item Erweiterte Matrix $[A\mid b]$ aufstellen.
    \item Mit elementaren Zeilenumformungen Zeilenstufenform erzeugen.
    \item Zuerst auf eine Widerspruchszeile prüfen.
    \item Bei einem lösbaren System freie Variablen wählen und die
      Pivotvariablen von unten nach oben bestimmen.
  \end{enumerate}
\end{KeyPointBox}

\begin{ExampleBox}{Vorwärtselimination und Rückwärtseinsetzen}{linear_systems_gaussian_example}
  \begin{ClozeArea}[34]{30}
    Ausgangssystem und erweiterte Matrix:
    \[
      \begin{aligned}
        2x_1+x_2-x_3&=3,\\
        3x_1+5x_2-4x_3&=1,\\
        4x_1-3x_2+2x_3&=2
      \end{aligned}
      \qquad\Longleftrightarrow\qquad
      \left[
        \begin{array}{rrr|r}
          2& 1&-1&3\\
          3& 5&-4&1\\
          4&-3& 2&2
        \end{array}
      \right].
    \]
    Die Operationen stehen jeweils rechts neben der Matrix, auf die sie
    angewendet werden:
    \[
      \begin{tikzpicture}[
  >=Latex,
  gaussmatrix/.style={
    matrix of math nodes,
    inner sep=0pt,
    left delimiter={[},right delimiter={]},
    nodes={minimum height=0.62cm,inner sep=0pt,anchor=center},
    column 1/.style={nodes={text width=0.42cm,align=right}},
    column 2/.style={nodes={text width=0.42cm,align=right}},
    column 3/.style={nodes={text width=0.42cm,align=right}},
    column 4/.style={nodes={text width=0.58cm,align=right}},
    column sep=0.35cm,row sep=1pt
  },
  operation/.style={font=\normalsize,color=black!55,anchor=west},
  operationbar/.style={draw=black!55,line width=0.6pt}
]
  \matrix (M0) [gaussmatrix] {
    2& 1&-1& 3\\
    3& 5&-4& 1\\
    4&-3& 2& 2\\
  };
  \draw[line width=0.6pt]
    ($(M0-1-3.north east)!0.5!(M0-1-4.north west)+(0,0.12)$)--
    ($(M0-3-3.south east)!0.5!(M0-3-4.south west)-(0,0.12)$);
  \coordinate (M0bar-top) at ($(M0.north east)+(0.75,0)$);
  \coordinate (M0bar-bottom) at ($(M0.south east)+(0.75,0)$);
  \coordinate (M0op) at ($(M0.east)+(0.98,0)$);
  \draw[operationbar] (M0bar-top)--(M0bar-bottom);
  \node[operation] at (M0op |- M0-2-4.center)
    {$2Z_2-3Z_1$};
  \node[operation] at (M0op |- M0-3-4.center)
    {$Z_3-2Z_1$};

  \matrix (M1) [gaussmatrix,below=0.85cm of M0] {
    2& 1&-1& 3\\
    |[text=THMBlue]| 0& 7&-5&-7\\
    |[text=THMBlue]| 0&-5& 4&-4\\
  };
  \draw[line width=0.6pt]
    ($(M1-1-3.north east)!0.5!(M1-1-4.north west)+(0,0.12)$)--
    ($(M1-3-3.south east)!0.5!(M1-3-4.south west)-(0,0.12)$);
  \coordinate (M1bar-top) at ($(M1.north east)+(0.75,0)$);
  \coordinate (M1bar-bottom) at ($(M1.south east)+(0.75,0)$);
  \coordinate (M1op) at ($(M1.east)+(0.98,0)$);
  \draw[operationbar] (M1bar-top)--(M1bar-bottom);
  \node[operation] at (M1op |- M1-3-4.center)
    {$7Z_3+5Z_2$};

  \matrix (M2) [gaussmatrix,below=0.85cm of M1] {
    2& 1&-1& 3\\
    |[text=THMBlue]| 0& 7&-5&-7\\
    |[text=THMBlue]| 0&|[text=THMBlue]| 0& 3&-63\\
  };

  \draw[line width=0.6pt]
    ($(M2-1-3.north east)!0.5!(M2-1-4.north west)+(0,0.12)$)--
    ($(M2-3-3.south east)!0.5!(M2-3-4.south west)-(0,0.12)$);
  \node[font=\normalsize,anchor=west] at ($(M2.east)+(0.98,0)$)
    {\Lang{Zeilenstufenform}{row echelon form}};
\end{tikzpicture}

    \]
    Die Koeffizientenmatrix besitzt drei Pivotpositionen, also Rang $3$.
    Nun wird von der letzten Gleichung nach oben eingesetzt:
    \begin{align*}
      3x_3&=-63
        &\Longrightarrow\qquad x_3&=-21,\\[0.7em]
      7x_2-5x_3&=-7
        &\Longrightarrow 7x_2-5(-21)&=-7,\\
      &&7x_2+105&=-7,\\
      &&7x_2&=-112,\\
      &&x_2&=-16,\\[0.7em]
      2x_1+x_2-x_3&=3
        &\Longrightarrow 2x_1+(-16)-(-21)&=3,\\
      &&2x_1+5&=3,\\
      &&2x_1&=-2,\\
      &&x_1&=-1.
    \end{align*}
    Somit lautet die eindeutige Lösung
    \[
      L=\{(-1,-16,-21)\}.
    \]
  \end{ClozeArea}
\end{ExampleBox}
